\section{Results}\label{sy-0200-ai}

The main results, with one proof deferred to the end of the section.

% !LOOM author: The synthetic quilt
% !LOOM created: 2026-09-16
% !LOOM tags: main

\begin{theorem}[Parity]\label{sy-0003-ai}
Let $(X,\sigma)$ be a finite widget. Then $|X| \equiv |\Fix(\sigma)| \pmod 2$.
\end{theorem}
\begin{proof}[First proof]\label{sy-0004-ai}
\uses{sy-0002-ai}
By Lemma~\ref{sy-0002-ai} the set $X$ is a disjoint union of orbits of size one or two.
\end{proof}
\begin{proof}[Second proof]\label{sy-0005-ai}
\uses{sy-0002-ai, sy-0001-ai}
Count fixed points of the permutation $\sigma$ directly, using \eqref{eq:step-ai} from the proof of Lemma~\ref{sy-0002-ai} and Definition~\ref{sy-0001-ai}.
\end{proof}


\subsection{Consequences}\label{sy-0201-ai}

\begin{proposition}[Two arguments]\label{sy-0006-ai}
A widget with $|X| = 1$ has a fixed point.
\end{proposition}
\begin{proof}
The only point is fixed.
\end{proof}
\begin{proof}
Apply Theorem~\ref{sy-0003-ai} with $|X|$ odd.
\end{proof}


% !LOOM tags: deferred

\begin{lemma}[Deferred]\label{sy-0007-ai}
The composite of two widget maps is a widget map.
\end{lemma}


\begin{definition}[Gadget]\label{sy-0008-ai}\label{def:gadget-ai}\label{defn:gadget-old-ai}
A \emph{gadget} is a widget whose involution has no fixed point.
\end{definition}


\begin{proof}[Proof of Lemma~\ref{sy-0007-ai}]
\uses{sy-0001-ai}
Both maps commute with the involutions, so their composite does. Here we use Definition~\ref{def:gadget-ai} only for notation.
\end{proof}



\subsection{Appendix: an incomplete argument}\label{sy-0300-ai}

This file is nested, so its section renders one level down.

\begin{lemma}[Incomplete]\label{sy-000C-ai}
Every gadget has even order.
\end{lemma}
\begin{proof}
\uses{sy-0008-ai}
Orbits have size two. \incomplete{Show that there are no fixed points using Definition~\ref{sy-0008-ai}.}
\end{proof}

\begin{lemma}[No proof yet]\label{sy-000D-ai}
Every widget map preserves orbit sizes.
\end{lemma}

\begin{remark}
An unlabelled remark: gadgets never have odd order.
\end{remark}

\begin{definition}[Half-gadget]\label{sy-000E-ai}
A widget with exactly one fixed point.
\end{definition}
\begin{proof}
Definitions need no proof; this one has a stray proof anyway.
\end{proof}

% !LOOM priority: high

\begin{proof}[Proof of Lemma~\ref{sy-000D-ai} and Lemma~\ref{sy-0007-ai}]
A proof naming two results attaches to the first.
\end{proof}


% A missing \input for the missing-include diagnostic. LaTeX would abort on it, so it sits inside \iffalse, which LaTeX skips and the scanner does not interpret.
\iffalse
\fi



\begin{mainthm}\label{sy-000B-ai}
\uses{sy-0003-ai, sy-0006-ai}
Every finite widget of odd order has a fixed point, and the fixed locus is functorial.
\end{mainthm}
\begin{proof}
Combine Theorem~\ref{sy-0003-ai} and Proposition~\ref{sy-0006-ai}.
\end{proof}


\bibliographystyle{amsplain}
\bibliography{refs}

